\section{Measurement Model}
\label{sec:measurement_model}

%%TODO: add corrected measurements for angular velocity and linear acceleration, as you referenced them before

We formulate the measurement models in both the \textit{right-invariant} and \textit{left-invariant} forms \textcolor{red}{(CITE this)}, i.e. $\mathbf{Y}_i = X^{-1} \mathbf{d}_i + \mathbf{V}_i$ for right-invariant and $\mathbf{Y}_i = X\mathbf{d}_i + \mathbf{V}_i$ for left-invariant. We define four measurement models, two in the right-invariant form (body frame centric), and two in the left-invariant form.

\subsection{Right-Invariant Measurement Models}

\subsubsection{Body-Frame Velocity Psuedo Measurement}

We define the body-frame velocity at the IMU origin as $\fv{B}{v} = \rotmat{W}{B}^\top \fv{W}{v}_{\mathcal{B}}$. And if we define $\mathbf{r}_w$ as the vector from the IMU origin to the center of the rear axle, by rigid body kinematics we have:

\begin{equation}
  \fv{B}{v}_{\mathrm{axle}} = \fv{B}{v}_{\mathrm{IMU}} + \bsomega^\wedge \mathbf{r}_w
  \label{eq:axle_vel}
\end{equation}

Now that we have the equation, we need to put this into Lie group measurement form. As this is body frame, we need to utilize the inverse of the state, as defined in \eqref{eq:inverse_state}. If we define the selection vector $\mathbf{d}=(\mathbf{0}_3,1,0)^\top$, then we get our measurement model as:

\begin{equation}
  X^{-1}\mathbf{d} =
  \begin{bmatrix}
    R^\top & -R^\top \mathbf{v} & -R^\top \mathbf{p} \\
    0 & 1 & 0 \\
    0 & 0 & 1
  \end{bmatrix}
  \begin{bmatrix}
    0 \\ 1 \\ 0
  \end{bmatrix}
  = \begin{bmatrix}
    -R^\top\mathbf{v} \\ 1 \\ 0
  \end{bmatrix}
  \label{eq:axle_meas}
\end{equation}

\subsubsection{Zero Angular Update Rate (ZARU)}

This measurement is deliberately excluded from the group, and is simply a method with which to pin $\mathbf{b}_g$ when the car is stationary, as at a zero point we know that $\tilde{\bsomega} = \boldsymbol{b}_g + \mathbf{n}_g$. Because of this, we can directly formulate the measurement model, along with its linearized Jacobian, as the following:

\begin{equation}
  \begin{gathered}
    \mathbf{z} = \tilde{\bsomega} - \hbs{b}_g \\
    H = \begin{bmatrix} 0 & 0 & 0 & -I_3 & 0 \end{bmatrix}
  \end{gathered}
  \label{eq:}
\end{equation}

\subsection{Left-Invariant Measurement Models}
Left-Invariant measurement model observes sensors attached to the robot that measure in world coordinates. 
\subsubsection{GPS Position Measurement}
We define the GPS position as the of the antenna in the ENU(East, North, Up) world frame. Let us define l as the vector going from the IMU to the GPS-antenna expressed relative to the body-frame. \begin{equation}
\begin{gathered}
        \fv{W}{\dbf{p}}_{antenna}=\fv{W}{\dbf{p}}_{\mathcal{B}}+\rotmat{W}{\mathcal{B}}\fv{\mathcal{B}}l
    \end{gathered}
\end{equation}
Defining our selection vector as $d=\begin{pmatrix}
    1^T,0,1
\end{pmatrix}^T$\begin{equation}
    Xd=\begin{bmatrix}
        R&v&p\\0&1&0\\0&0&1
    \end{bmatrix}\begin{bmatrix}
        1\\0\\1
    \end{bmatrix}=\begin{bmatrix}
        Rl+p\\0\\1
    \end{bmatrix}
\end{equation}
